\begin{resumo}[Abstract]
\begin{otherlanguage*}{english}
Coastal biodiversity monitoring is essential for environmental conservation but faces challenges related to data collection and quality. This work, in collaboration with the Center for Coastal, Limnological and Marine Studies (CECLIMAR), aimed to develop a mobile application called \textit{FaunaMar}, designed to standardize and improve the efficiency of the data lifecycle in monitoring coastal fauna along the northern coast of Rio Grande do Sul, Brazil.
The project methodology combined a bibliographic and documentary research phase to support the theoretical foundation with an agile development process, managed using Kanban via the \textit{Jira} platform. Requirements were gathered through meetings with CECLIMAR professionals, and the interface was prototyped using \textit{Figma}. The application was developed using \textit{Flutter} and the \textit{Firebase} platform as the backend, which included services such as \textit{Cloud Firestore}, \textit{Authentication}, and \textit{Storage}, along with the \textit{Hive} database to ensure offline functionality.
The results include a functional application that digitizes and centralizes the monitoring process. The system offers two user profiles: citizen scientist and researcher, each with distinct interfaces and permissions. Key features include simple and technical registration forms, GPS-based geolocation capture, and a record dashboard with advanced filters and an interactive map. Additionally, it includes a gamified user profile with achievements and a dedicated module for researchers to evaluate submissions. The project followed a structured and traceable development process with 131 tasks, 86\% of which were completed. The application was validated through 17 test versions with 7 users from different profiles, confirming its usability. The delivered technological solution meets a real demand from CECLIMAR, applying Citizen Science as a tool for environmental research. The FaunaMar app enhances the quality and efficiency of coastal fauna monitoring, reduces researchers' operational workload, and strengthens the connection between science and society in promoting a more sustainable future.

\textbf{Keywords}: citizen science. environmental monitoring. mobile application.
\end{otherlanguage*}
\end{resumo}